%% ==========================================================================
%% working_paper.tex — 正文工作文件
%% 首投：Automation in Construction (Elsevier)，模板 elsarticle
%% 结构说明见 paper/PAPER_PLAN.md；参考文献 paper/references.bib
%% 编译：latexmk -pdf working_paper.tex（或 pdflatex -> bibtex -> pdflatex x2）
%% 首投可用 "Your Paper Your Way"（本模板即够）；转投 MDPI 时换模板搬正文
%% 红线：正文数字一律核 phase6_gen493_summary.json（复核修订版）/ PROGRESS.md 三（结果与结论），
%%       不在此文件手抄改数（防口径漂移）
%% ==========================================================================
\documentclass[preprint,12pt]{elsarticle}

\usepackage[T1]{fontenc}
\usepackage[utf8]{inputenc}
\usepackage{graphicx}
\usepackage{booktabs}
\usepackage{amsmath,amssymb}
\usepackage{siunitx}
\usepackage{url}
\usepackage[hidelinks]{hyperref}
\usepackage{xcolor}

% 占位标记：红字 TODO，定稿前全文清零
\newcommand{\todo}[1]{\textcolor{red}{[TODO: #1]}}

\graphicspath{{figures/}}

\journal{Automation in Construction}

\begin{document}

\begin{frontmatter}

\title{\todo{Title:\\Generative data augmentation for scarce hazardous targets in construction safety: a reviewed pipeline and quantity-over-quality evidence (working title)}}

\author[1]{\todo{Author names}}
\affiliation[1]{organization={\todo{Institution}}, city={\todo{City}}, country={\todo{Country}}}

\begin{abstract}
\todo{Abstract (<=200 words): problem (scarce dangerous samples, can't stage violations safely) ->
pipeline (diffusion generation + auto-labeling + human review) ->
evidence (YOLO detectors trained on generated data outperform real-only baseline on clean real test;
quantity-over-quality scaling) -> implication (reusable for other scarce targets).}
\end{abstract}

\begin{keyword}
synthetic data \sep data augmentation \sep construction safety \sep
object detection \sep diffusion models \sep human-in-the-loop annotation
\end{keyword}

\end{frontmatter}

% ==========================================================================
\section{Introduction}\label{sec:intro}
% 写作要点（PAPER_PLAN.md 第二节）：
%  1) 危险样本稀缺是趋势性的（监管趋严 -> 违规样本减少），但后果严重必须检出
%  2) 摆拍违规场景有安全/伦理问题，不可行
%  3) 生成式扩量是出路；真实性是红线（v0 教训：过净合成图 -> 域差）
%  4) 本文：完整管线 + 干净评测协议 + 数量取胜实证 +（第二目标后）跨目标通用性
% 贡献列表（等 4.4/4.6 实验完成后定稿；当前占位）：
%  C1: an end-to-end pipeline (generation -> auto-labeling -> human-in-the-loop review)
%      for synthesizing scarce hazardous-scenario data, demonstrated on gas cylinders
%      with placement issues;
%  C2: a leakage-aware evaluation protocol that keeps all real photos out of training
%      and quantifies the inflation caused by conventional split leakage;
%  C3: empirical evidence that increasing the *quantity* of diverse generated samples
%      improves real-world detection (scaling ablation, \S\ref{sec:scale});
%  C4 (optional): generality to a second scarce target (helmet non-use, \S\ref{sec:helmet}).
\todo{英文初稿（依赖文献阅读；引 \cite{nath2020ppe,jeong2026lora,davari2025bcon} 作动机支撑）}

% ==========================================================================
\section{Related Work}\label{sec:related}

\subsection{Computer vision for construction safety}
% 已有本地文献：\cite{nath2020ppe} PPE 实时检测；\cite{fang2018nocert} 非持证作业；
% \cite{zhang2025helmet} 安全帽；数据集 \cite{duan2022soda}。
% 综述线索：工地 CV 长期受限于数据获取（安全/隐私/成本）。
\todo{编写：从 PPE/行为检测讲到数据瓶颈}

\subsection{Synthetic and generative data for object detection}
% 三条线：
%  a) 图形引擎 sim2real + 域随机化：\cite{richter2016playing, tremblay2018domainrand, lee2023gameengine}
%  b) 生成模型扩增：扩散基础 \cite{ho2020ddpm, rombach2022ldm, zhang2023controlnet}；
%     下游增广 \cite{azizi2023synthetic, ghiasi2021copypaste, shorten2019augmentation}
%  c) 建筑/工业场景就近工作（重点差异化对象）：
%     \cite{davari2025bcon}（BlendCon+ControlNet，worker 检测，25.6k 图）
%     \cite{jeong2026lora}（LoRA 微调，危险场景 falls/struck-by/caught-in）
%     \cite{zhang2026mixing}（真实-合成混合固定预算）
%     \cite{girella2024diag, li2025steeldefect, liu2025substation, dinh2026ponding}
% 差异点：本文针对"违规状态"类危险样本（不是 worker 本体），且强调
%        全自动打标+人工复核闭环 与 干净评测协议+数量缩放实证。
\todo{编写（这是与最接近竞品划界限的关键小节）}

\subsection{Auto-labeling and human-in-the-loop annotation}
% \cite{lee2013pseudolabel, xie2020noisystudent, northcutt2021labelerrors}
% 本文闭环：打标模型预标注 -> GUI 人工复核（复核率 100%），人工复核是质量门。
\todo{编写}

% ==========================================================================
\section{Methodology}\label{sec:method}

\subsection{Target definition and problem statement}
% "Placement Issues" 口径（监理实际管理口径）：气瓶未放置在专用固定架/手推车上
% 即为正样本（无论倒地与否）；规范放置=负样本（背景图）。
\todo{编写（素材：PROGRESS.md 一（项目与目标））}

\subsection{Real seed data and annotation}
% 93 张真实巡检照片：57 正（179 框）+ 36 负（空标），人工标注；
% 数据来源与合规说明。
\todo{编写（素材：PROGRESS.md 二.1）}

\subsection{State-editing synthesis pipeline}
% Qwen-Image-Edit-2509 图生图编辑真实巡检照片 -> 合成"违规状态"样本
% （两批共 1085 张，全部源自 93 张真实照片；逐图来源可审计）。
% 真实性红线：必须带喷漆磨损/污渍/锈迹（v0 教训）；来源记录 = 方法学贡献之一。
\todo{编写（素材：PROGRESS.md 一/二；工具引用 \cite{qwenimage2025, qwenimageedit2025}）}

\subsection{Auto-labeling with human-in-the-loop review}
% 打标模型（gen1085/yolo26s，在 1085 张合成图上训练）预标注 ->
% annotator GUI 逐张人工复核（100% 复核率）-> 已复核合成数据（1085 张）。
\todo{编写（素材：PROGRESS.md 二.2 与 五（打标模型））}

\subsection{Evaluation protocol}\label{sec:protocol}
% 关键设计（本文方法学贡献之一）：
%  - test = 61 张网络来源图（72 框），与所有训练数据来源完全独立；
%  - 合成数据的来源审计协议：逐图记录编辑源，test 图片不得作为生成源
%    （v2 曾因此泄漏 611/1085 张，公开该教训与审计工具）；
%  - 全篇统一：固定 100 epoch / imgsz 640 / batch 16、报 last.pt、不做选权；
%  - 三臂同源对照（同一批 93 张真实照片：原样 / 传统增广 / 编辑合成）。
\todo{编写（素材：PROGRESS.md 三（设置与结果）；口径见 List_of_Experiments.md）}

\begin{figure}[t]
  \centering
  \fbox{\parbox{0.92\linewidth}{\centering
    \todo{Fig.~1 管线总图\\（先用 \texttt{paper/fig1\_pipeline.md} 的 mermaid 草图记录思路；\\
    定稿前用 TikZ 正式重绘并替换 \texttt{\textbackslash includegraphics\{fig1\_pipeline.pdf\}}）}}}
  \caption{\todo{Overview of the proposed pipeline: real-photo seeding, generative
  expansion, auto-labeling with human-in-the-loop review, and leakage-aware
  training/evaluation; the loop re-instantiates for other scarce hazardous targets.}}
  \label{fig:pipeline}
\end{figure}

% ==========================================================================
\section{Experiments and Results}\label{sec:exp}

\subsection{Experimental setup}
% 6 个检测器（YOLOv8/v11/v26 x s/m），统一超参 epochs=100 / imgsz=640 /
% batch=16 / device=0，报最后一轮（last.pt，不做验证集选权）；
% 硬件 RTX 5070 Ti Laptop (12 GB)；各类计数。
\todo{编写（素材：PROGRESS.md 三（结果）；引用 \cite{redmon2016yolo,
bochkovskiy2020yolov4, wang2023yolov7, wang2024yolov9, wang2024yolov10, ultralytics2025}）}

\subsection{Main comparison: three same-source arms}
% A 真实 93 张 / B 传统增广 868 / C 编辑合成 868，全部只差训练数据来源，
% 统一在新 test 上评测（六权重）。C 臂已出（0.768~0.927 mAP50）；A/B 待训。
\todo{编写 + 表（素材：phase10_v3_main.json + PROGRESS.md 三）}

\subsection{Effect of generated data on real-world detection}\label{sec:main}
% 主结果表：三臂 × 六权重 × 新 test（A/B 训完后汇总）。
% 当前已有：C 臂六权重 + gen493 规模点。
\begin{table}[t]
  \centering
  \caption{\todo{Detection performance on the independent test set (61 web images).}}
  \begin{tabular}{lcccc}
  \toprule
  Detector & \multicolumn{2}{c}{Real-only (A)} & \multicolumn{2}{c}{Synthesis (C)} \\
  \cmidrule(lr){2-3}\cmidrule(lr){4-5}
   & mAP50 & mAP50--95 & mAP50 & mAP50--95 \\
  \midrule
  \todo{6 行数据} & & & & \\
  \bottomrule
  \end{tabular}
  \label{tab:main}
\end{table}

\subsection{Data-scale ablation}\label{sec:scale}
% ✅ 已完成（2026-09-19）：gen493（train 394）vs gen1085（train 868）在新 test 上，
% 六权重 mAP50-95 全涨 +0.07~+0.15（yolo26s 0.557→0.698）。
\todo{编写：规模曲线 + 小表（素材：phase10_v3_main.json）}

\subsection{Inference resolution analysis}
% imgsz 扫描：640/960/1280/1536；640 峰值；真实照片 17% 框在 640 等效下 <20px。
% 表 + 讨论小目标因素（\cite{cheng2023smallobject, akyon2022sahi}）。
\todo{编写（素材：PROGRESS.md 三（imgsz 结论④））}

\subsection{Cross-target generalization: helmet non-use}\label{sec:helmet}
\todo{待实验：安全帽第二目标（管线复用）。三级定位第 2 级的门槛，见 PAPER_PLAN 第五节。}

\subsection{Comparison with classical augmentation and public datasets}
\todo{待实验：传统增广臂（B）= 93 张离线增广到 868（与 C 臂对齐张数与正负比），
统一协议与新 test（\cite{ghiasi2021copypaste, shorten2019augmentation}）；
设计见 PAPER_PLAN 第四节 1。公开数据：两个 Roboflow 气瓶数据集无放置位姿标注、
不可用（\cite{rfcylinders, rfgarrafas}），改为 Introduction/Discussion 论据。}

% ==========================================================================
\section{Discussion}\label{sec:discussion}
% 要点：
%  1) domain gap 两个来源：合成图"过净"（v0 教训）+ 小目标被分辨率压缩（本工作发现）
%  2) 评测协议启示：数据扩增类论文的泄露陷阱（\cite{kaufman2012leakage, northcutt2021labelerrors}）
%  3) 伦理与合规：生成式 AI 使用披露、生成数据不用于冒充真实证据；标注人工复核质量门
%  4) 通用性愿景（勿 overclaim）：can be re-instantiated for other scarce targets;
%     跨行业图像识别数据稀缺皆适用（Discussion/future work 口径，见 PAPER_PLAN 第五节三级定位）
\todo{编写}

% ==========================================================================
\section{Conclusions}\label{sec:conclusion}
\todo{最后写：3~5 句 + 局限 + future work}

% ==========================================================================
% AiC 投稿要求的声明节（占位，投稿前完善）
\section*{CRediT authorship contribution statement}
\todo{roles: conceptualization / methodology / software / ...}

\section*{Declaration of competing interest}
The authors declare that they have no known competing financial interests or
personal relationships that could have appeared to influence the work reported
in this paper.

\section*{Data availability}
\todo{state what can be shared: annotation protocols, splits, model weights;
generated datasets may be released upon request / with the paper.}

% ==========================================================================
\bibliographystyle{elsarticle-num}
\bibliography{references}

\end{document}
